\documentclass[11pt]{article}
\usepackage[margin=1in]{geometry}
\usepackage[T1]{fontenc}
\usepackage[utf8]{inputenc}
\usepackage{amsmath,amssymb,amsthm}
\usepackage{enumitem}

\title{ICPC World Finals 2020\\L. Sweep Stakes}
\author{}
\date{}

\begin{document}
\maketitle

\section*{Problem Summary}

Each cell $(i,j)$ independently contains a mine with probability $a_{ij}=p_i+q_j$.
We are given that the total number of mines is exactly $t$.
For each query subset $S$ of at most $500$ cells, we must output the conditional distribution of the number of mines in
$S$.

\section*{Key Observations}

\begin{itemize}[leftmargin=*]
    \item Let $X=\sum X_{ij}$ be the total number of mines.
    The hard part is the conditioning on $X=t$.
    A standard exponential tilt removes that difficulty.
    \item For any $z>0$, define a new independent model by
    \[
        \Pr_z(X_{ij}=1)=b_{ij}(z)=\frac{a_{ij}z}{1-a_{ij}+a_{ij}z}.
    \]
    For every full configuration $\omega$ with exactly $t$ mines,
    \[
        \Pr_z(\omega)=C(z)\,z^t\,\Pr(\omega),
    \]
    where $C(z)$ depends only on $z$, not on $\omega$.
    Therefore the conditional distribution given $X=t$ is \emph{identical} in the original and tilted models.
    \item We choose $z$ so that $\mathbb E_z[X]=t$.
    Then the event $X=t$ becomes a central event in the tilted model.
    \item For a query subset $S$, let $Y$ be the number of mines inside $S$ and let $Z$ be the number of mines outside
    $S$.
    In the tilted model, $Y$ and $Z$ are independent, so
    \[
        \Pr(Y=r\mid X=t)=\frac{\Pr_z(Y=r)\Pr_z(Z=t-r)}{\Pr_z(X=t)}.
    \]
    The inside part $\Pr_z(Y=r)$ is exact by a tiny $O(s^2)$ DP because $s\le 500$.
    \item We only need point probabilities near the mean.
    So instead of storing the whole distribution of $X$, we store the distribution of $X \bmod M$ for a carefully
    chosen odd modulus $M$.
    If $M$ is large enough, the probability mass at the other congruent integers is exponentially tiny.
\end{itemize}

\section*{Algorithm}

\subsection*{1. Choose the tilt}

Let
\[
    f(z)=\sum_{i,j}\frac{a_{ij}z}{1-a_{ij}+a_{ij}z}.
\]
This is strictly increasing from $0$ to $mn$, so binary search finds the unique $z$ with $f(z)=t$.
Then compute all tilted probabilities
\[
    b_{ij}= \frac{a_{ij}z}{1-a_{ij}+a_{ij}z}.
\]

\subsection*{2. Choose a modulus}

Let
\[
    \sigma^2=\sum_{i,j} b_{ij}(1-b_{ij})
\]
be the tilted variance of the total count.
For the largest query size $s_{\max}$, choose the smallest odd $M>s_{\max}$ such that
\[
    2\exp\!\left(-\frac{(M-s_{\max})^2}{2(\sigma^2+(M-s_{\max})/3)}\right) < 10^{-11}.
\]

This is Bernstein's tail bound. It ensures that any count at distance at least $M-s_{\max}$ from the mean has
negligible probability, which is enough for the final $10^{-6}$ error bound.

\subsection*{3. Build the total distribution modulo $M$}

Let $D[r]=\Pr_z(X\equiv r \pmod M)$.
We build it by a cyclic DP over all cells.
If the next cell has probability $p$, then
\[
    D_{\text{new}}[r]=(1-p)D[r]+pD[r-1]
\]
with indices modulo $M$.
This costs $O(mnM)$.

\subsection*{4. Answer one query}

Suppose the query cells are $c_1,\dots,c_s$.

\paragraph{Inside part.}
Run the ordinary $O(s^2)$ Poisson-binomial DP on the probabilities $b_{c_1},\dots,b_{c_s}$ to get
\[
    I[r]=\Pr_z(Y=r).
\]

\paragraph{Outside part modulo $M$.}
Start from the already computed total residue distribution $D$ and remove the selected cells one by one.

For one cell with probability $p$, if $C$ is the residue distribution before adding that cell and $D$ is the residue
distribution after adding it, then
\[
    D[r]=(1-p)C[r]+pC[r-1].
\]
So answering the query requires inverting this linear transform $s$ times.

Because $M$ is odd, the transform is invertible for every $0<p<1$:
its eigenvalues are $1-p+p\omega^j$ for the $M$-th roots of unity $\omega^j$, and the only way such a value could be
zero is $p=\tfrac12$ and $\omega^j=-1$, impossible for odd $M$.

The inversion itself is done in $O(M)$ by a simple recurrence:
\begin{itemize}[leftmargin=*]
    \item if $p\le \tfrac12$, solve forward using
    \[
        C[r]=\frac{D[r]-pC[r-1]}{1-p};
    \]
    \item if $p> \tfrac12$, solve backward using
    \[
        C[r-1]=\frac{D[r]-(1-p)C[r]}{p}.
    \]
\end{itemize}
In either direction one unknown initial value remains; it is determined by the final wrap-around equation.

After removing all selected cells we obtain
\[
    O[r]=\Pr_z(Z\equiv r \pmod M).
\]

\paragraph{Final formula.}
For every $r=0,\dots,s$ output
\[
    \frac{I[r]\cdot O[(t-r)\bmod M]}{D[t\bmod M]}.
\]
Because of the chosen bound on $M$, the residue probabilities differ from the true point probabilities by far less than
the allowed error.

\section*{Correctness Proof}

We prove that the algorithm outputs the correct probabilities.

\paragraph{Lemma 1.}
For any $z>0$, conditioning on the total number of mines being $t$ gives the same distribution in the original and
tilted models.

\paragraph{Proof.}
Fix a full configuration $\omega$ and let $|\omega|$ be its number of mines.
In the original model,
\[
    \Pr(\omega)=\prod_{\omega_{ij}=1} a_{ij}\prod_{\omega_{ij}=0}(1-a_{ij}).
\]
In the tilted model,
\[
    \Pr_z(\omega)=\prod_{\omega_{ij}=1}\frac{a_{ij}z}{1-a_{ij}+a_{ij}z}
    \prod_{\omega_{ij}=0}\frac{1-a_{ij}}{1-a_{ij}+a_{ij}z}.
\]
Therefore
\[
    \Pr_z(\omega)=
    \left(\prod_{i,j}\frac{1}{1-a_{ij}+a_{ij}z}\right)
    z^{|\omega|}\Pr(\omega).
\]
If $|\omega|=t$, the multiplicative factor depends only on $z$ and $t$, not on $\omega$.
So after normalizing over all configurations with total $t$, both conditional distributions are identical. \qed

\paragraph{Lemma 2.}
For any query subset $S$ and any $r$,
\[
    \Pr(Y=r\mid X=t)=\frac{\Pr_z(Y=r)\Pr_z(Z=t-r)}{\Pr_z(X=t)}.
\]

\paragraph{Proof.}
By Lemma~1 we may work in the tilted model.
There the cells are independent, so the counts inside and outside the query are independent:
\[
    \Pr_z(Y=r, Z=t-r)=\Pr_z(Y=r)\Pr_z(Z=t-r).
\]
Since $X=Y+Z$, conditioning on $X=t$ gives exactly the stated formula. \qed

\paragraph{Lemma 3.}
The one-cell residue update
\[
    D[r]=(1-p)C[r]+pC[r-1]
\]
is invertible for odd $M$ and $0<p<1$.

\paragraph{Proof.}
This is a circulant linear map on residue vectors of length $M$.
Its eigenvalues in the Fourier basis are
\[
    \lambda_j=1-p+p\omega^j,
\]
where $\omega$ is a primitive $M$-th root of unity.
If $\lambda_j=0$, then $p\omega^j=-(1-p)$.
Taking absolute values gives $p=1-p$, so $p=\tfrac12$, and then $\omega^j=-1$.
But $-1$ is not an odd-order root of unity.
Hence no eigenvalue is zero, so the map is invertible. \qed

\paragraph{Lemma 4.}
For each query, the algorithm computes the exact residue distribution of the outside count $Z \bmod M$.

\paragraph{Proof.}
The total residue distribution is obtained by repeatedly applying the one-cell forward update, so it is exact.
Removing one selected cell applies the inverse of that exact linear map, which exists by Lemma~3.
After removing all selected cells, what remains is exactly the distribution of the sum of all unselected cells modulo
$M$, i.e.\ the residue distribution of $Z$. \qed

\paragraph{Lemma 5.}
For every quantity used in the final answer, replacing a point probability by the corresponding residue probability
changes the value by at most the Bernstein tail bound used when choosing $M$.

\paragraph{Proof.}
In the tilted model the total count has mean $t$.
For a query of size $s$, the outside count $Z$ has mean $t-\mathbb E_z[Y]$, so the requested point $t-r$ differs from
that mean by at most $s$.
Any other integer congruent to $t-r$ modulo $M$ is therefore at least $M-s$ away from the mean.
The residue probability is the sum of the desired point probability and all those other congruent point probabilities,
so the aliasing error is bounded by the tail probability outside distance $M-s$.
Bernstein's inequality gives exactly the bound used by the algorithm. \qed

\paragraph{Theorem.}
The algorithm outputs every query answer with the required accuracy.

\paragraph{Proof.}
By Lemma~2 the exact answer factorizes into the inside point probability, the outside point probability, and the total
point probability.
The inside part is computed exactly by the small DP.
By Lemma~4 the outside and total residue distributions are exact modulo $M$, and by Lemma~5 the replacement of point
probabilities by these residue probabilities introduces only a negligible error.
Therefore every reported probability is within the required tolerance. \qed

\section*{Complexity Analysis}

Let $N=mn$ and let $M$ be the chosen odd modulus.

\begin{itemize}[leftmargin=*]
    \item Solving for the tilt takes $O(N \log \text{precision})$.
    \item Building the total residue distribution takes $O(NM)$.
    \item A query of size $s$ takes $O(s^2+sM)$.
\end{itemize}

So the whole program runs in
\[
    O\!\left(NM + \sum_{\text{queries}} (s^2+sM)\right)
\]
time and
\[
    O(N+M)
\]
memory.

With the given constraints, the chosen $M$ stays small (about a few thousand at worst), so this easily fits.

\section*{Implementation Notes}

\begin{itemize}[leftmargin=*]
    \item The cases $t=0$ and $t=mn$ are deterministic and are handled separately.
    \item The forward inversion is numerically stable when $p\le \tfrac12$, and the backward inversion is numerically
    stable when $p>\tfrac12$.
    \item After each query, the output vector is renormalized to sum to $1$ in order to remove tiny floating-point
    drift.
\end{itemize}

\end{document}
